\begin{textsample}{NEG Dim 4 – 2018\_c6 – Score 12.00 – 2018\_c6/enade\_2018\_secr\_exec\_32b5ea331f.txt}  \label{ex:f4_neg_339}
CAPÍTULO 6 ANÁLISE TÉCNICA DA PROVA Este capítulo tem por objetivo apresentar o desempenho dos estudantes concluintes de Secretariado Executivo no Enade/2018.
Para isso, foram calculadas as estatísticas básicas da prova em seu todo ( seção 6.1.1 ), bem como as estatísticas relacionadas aos Componentes de Formação Geral ( seção 6.1.2 ) e de Conhecimento Específico da Área ( seção 6.1.3 ).
Dadas as suas características, foram analisadas, em separado, as questões objetivas ( seção 6.2 ) e as questões discursivas ( seção 6.3 ).
Para as questões objetivas, foram disponibilizados os índices de facilidade e de discriminação Ponto-Bisserial, também em separado, para os Componentes de Formação Geral ( seção 6.2.1 ) e de Conhecimento Específico ( seção 6.2.2 ).
De cada componente, uma das questões foi escolhida para exemplificar a análise gráfica, relacionando as alternativas escolhidas pelos estudantes ( inclusive o gabarito ) com o número de acertos no componente.
No Anexo I, apresenta -se a íntegra da análise gráfica para todas as questões objetivas.
Para cada uma das questões discursivas, os conteúdos dos tipos mais comuns de respostas dos estudantes são apresentados e comparados com o padrão de respostas esperado ( ver Anexo VIII com o padrão de respostas ).
Tomando -se como base as duas questões discursivas do Componente de Formação Geral, na seção 6.3.1.6, são apresentados comentários sobre a correção das respostas em relação à Língua Portuguesa.
Nas tabelas, constam as seguintes estatísticas das notas²⁰ : média do desempenho na prova, erro padrão da média, desvio padrão, nota mínima, mediana e nota máxima para cada um de seus componentes.
Tais estatísticas contemplam o total de estudantes concluintes da Área de Secretariado Executivo inscritos e presentes à prova do Enade/2018, tendo em vista agregações, ou por Grandes Regiões ou o país como um todo, ou por Categoria Administrativa, Organização Acadêmica da IES e Modalidade de Ensino.
Em relação aos gráficos de distribuição de notas, o intervalo considerado foi de 10 unidades, aberto à esquerda e fechado à direita, com exceção do primeiro intervalo, [ 0 ; 10 ], fechado em ambos os extremos.
Para os gráficos de distribuição das notas das questões discursivas, foram consideradas mais duas categorias : questão em branco²¹ e nota zero. ²⁰ Uma definição dessas estatísticas pode ser encontrada no Glossário. ²¹ Nesse grupo estão incluídas também as respostas classificadas como nulas ou desconsideradas. 6.1 ESTATÍSTICAS BÁSICAS DA PROVA Esta seção apresenta estatísticas selecionadas e histogramas da nota geral ( 6.1.1 ) e de cada componente : Formação Geral ( 6.1.2 ) e Conhecimento Específico ( 6.1.3 ).
São também, apresentadas estatísticas selecionadas de subpopulações, caracterizadas por Grande Região, Categoria Administrativa, Organização Acadêmica e Modalidade de Ensino. 6.1.1 Estatísticas Básicas Gerais Na Tabela 6.1, são apresentadas as Estatísticas Básicas da prova, por Grande Região, dos estudantes concluintes de Secretariado Executivo.
A Média das notas da prova, como um todo ( nas seções seguintes serão analisados os Componentes de Formação Geral e de Conhecimento Específico ), foi 52,3, sendo que os alunos da região Centro-Oeste obtiveram a Média mais baixa ( 46,0 ), e os da região Sudeste obtiveram a Média mais alta ( 55,3 ).
As demais médias foram : 46,3, na região Norte, 53,4, na região Nordeste e 54,1, na região Sul.
O Desvio padrão para o Brasil, como um todo, foi 14,1, sendo o maior Desvio padrão encontrado na região Centro-Oeste ( 15,2 ), e o menor, na região Norte ( 11,5 ), indicando uma dispersão um pouco menor das notas dessa última região.
A região que obteve a maior nota Máxima foi a Sudeste ( 89,0 ), ao passo que a região que atingiu a menor nota Máxima foi a Norte ( 73,6 ).
A Mediana do Brasil, como um todo, foi 52,6, sendo a maior Mediana obtida na região Sudeste ( 55,6 ), e a menor obtida na Centro-Oeste ( 44,5 ).
A nota Mínima do Brasil foi obtida na região Centro-Oeste, onde esta foi zero.
A Mínima na região Norte foi 24,4, a maior nota Mínima dentre as cinco Grandes Regiões, e, nas demais, foram : 17,6, na região Nordeste, 16,5, na Sudeste e 17,1, na Sul.
Considerando -se as notas segundo Grande Região, observa -se que não existe diferença estatisticamente significativa ao nível de 95 \% entre as duas menores Médias, obtidas nas regiões Centro-Oeste e Norte, e nem entre as três maiores, nas regiões Nordeste, Sudeste e Sul²². ²² Todas as comparações deste capítulo utilizam os intervalos de 95 \%.
Os erros-padrão da média que possibilitam os testes estão disponíveis nas tabelas.
O comportamento das notas dos estudantes de todo o Brasil pode ser observado no Gráfico 6.1, que apresenta um histograma com a distribuição das mesmas : uma distribuição unimodal com a moda no intervalo ( 50,60 ].
Na Tabela 6.2, são apresentadas informações referentes à Média da nota final, desagregadas por Categoria Administrativa, por Organização Acadêmica e por Modalidade de Ensino.
A Média das notas da prova, como um todo, foi 52,3.
Em relação à Categoria Administrativa, os estudantes das IES Públicas obtiveram Média mais alta ( 52,8 ) que a Média No tocante à Organização Acadêmica, as Universidades e os Centros Universitários obtiveram Média mais alta que a nacional ( 53,7 e 55,1, respectivamente ).
A Média das Faculdades e a dos CEFET/IFET foram menores do que a nacional ( 46,2 e 46,5, respectivamente ).
Constata -se que entre as Organizações Acadêmicas existe diferença estatisticamente significativa ao nível de 95 \% entre as duas Médias mais altas e as duas mais baixas.
A Média da Modalidade Educação a Distância ( 51,2 ) foi inferior à Média da Educação Presencial ( 52,3 ), e a diferença entre elas não é estatisticamente significativa. 6.1.2 Estatísticas Básicas no Componente de Formação Geral Na Tabela 6.3, são apresentadas as Estatísticas Básicas em relação ao componente da prova que avalia a Formação Geral dos estudantes concluintes.
Os alunos de todo o Brasil obtiveram desempenho médio de 42,5.
Quanto à variabilidade, o Desvio padrão das notas dos estudantes do Brasil, como um todo, foi 16,9.
A maior Média foi obtida na região Sudeste A maior nota no Componente de Formação Geral da prova do Enade/2018 foi 95,6, obtida por, pelo menos, um aluno na região Nordeste.
A menor nota Máxima foi obtida na região Centro-Oeste ( 75,9 ).
A Mediana do Brasil, como um todo, foi 41,5, sendo a menor Mediana encontrada na região Centro-Oeste ( 33,1 ), e a maior, na região Sul ( 44,6 ).
A menor nota Mínima nessa parte foi zero, nas regiões Sul e Centro-Oeste, e a maior nota Mínima foi 19,5, na Norte.
Considerando -se as notas, segundo Grande Região, observa -se que só existe diferença estatisticamente significativa entre as Médias das notas do Componente de Formação Geral obtidas na região Centro-Oeste e as regiões Nordeste, Sudeste e Sul.
No Gráfico 6.2, é possível verificar a avaliação do desempenho dos estudantes do Componente de Formação Geral, a partir do histograma da distribuição das notas correspondentes.
A distribuição é unimodal, com moda em ( 30 ; 40 ], seguida de perto pelo intervalo ( 40 ; 50 ].
Nota -se, ainda, que, no Gráfico 6.2, as notas apresentam maior dispersão do que as no Gráfico 6.1 ( distribuição das notas da prova ), confirmado pela comparação dos desvios padrão : 14,1 para a nota da prova, como um todo, e 16,9 para o Componente de Formação Geral.
Na Tabela 6.4, são apresentadas as informações referentes ao desempenho dos concluintes do Componente de Formação Geral, em diferentes agregações : Categoria Administrativa, Organização Acadêmica e Modalidade de Ensino.
Observa -se que não existe diferença estatisticamente significativa entre as médias dos tipos de Categoria Administrativa : a maior Média obtida por estudantes de IES Públicas ( 43,5 ) e a menor, pelos de IES Privadas ( 41,1 ).
Nota -se, também, que há diferença estatisticamente significativa entre a maior das médias dos quatro tipos de Organização Acadêmica e a menor delas, o valor maior para as Universidades ( 44,3 ) e o menor para os CEFET/IFET ( 35,9 ).
No que se refere à Modalidade de Ensino, a Educação Presencial apresentou maior Média do que a Educação a Distância ( médias 42,5 e 37,0, respectivamente ).
Tais médias não apresentam, entretanto, diferença estatisticamente significativa. 6.1.3 - Estatísticas Básicas do Componente de Conhecimento Específico Na Tabela 6.5, são apresentadas as Estatísticas Básicas referentes ao Componente de Conhecimento Específico da área de Secretariado Executivo.
A Média do desempenho dos alunos do Brasil, como um todo, foi 55,5.
A maior Média foi obtida na região Sudeste ( 58,7 ), e a menor, na região Norte ( 48,7 ).
As demais médias foram : 56,4, na região Nordeste, 57,3, na região Sul e 49,3, na região Centro-Oeste.
Quanto à variabilidade das notas, o Desvio padrão do Brasil, como um todo, foi 15,1, sendo o maior Desvio padrão observado na região Centro-Oeste ( 16,7 ), e o menor, nas regiões Norte e Sudeste ( 13,3 ).
Os demais desvios foram : 15,0, na região Nordeste e 14,1, na região Sul.
A Mediana das notas dos estudantes de todo o Brasil foi 56,8.
A maior Mediana ocorreu na região Sudeste ( 59,4 ), e a menor, na região Norte ( 48,5 ).
As demais medianas foram : 57,5, na região Nordeste, 57,0, na região Sul e 49,1, na região Centro-Oeste.
A nota Máxima do Brasil, como um todo, foi 90,6, sendo obtida por, pelo menos, um aluno na região Nordeste.
As demais notas máximas foram : 72,6, na região Norte, 87,9, na região Sudeste, 87,7, na região Sul e 81,7, na região Centro-Oeste.
A nota Mínima foi zero, obtida na região Centro-Oeste.
Nas demais regiões, as notas mínimas foram : 23,2, na região Norte, 14,4, na região Nordeste, 10,7, na região Sudeste e 13,1, na região Sul.
Observa -se que não existe diferença estatisticamente significativa entre as duas menores Médias, das notas do Componente de Conhecimento Específico, das regiões Norte e Centro-Oeste.
No entanto, existe diferença estatisticamente significativa entre essas duas regiões, que obtiveram as menores médias, e as demais.
Assim como nos Gráficos 6.1 e 6.2, no Gráfico 6.3, apresentado a seguir, é possível ser feita uma avaliação do desempenho de concluintes de Secretariado Executivo, em relação ao Componente de Conhecimento Específico, com um histograma da distribuição das notas correspondentes.
Esta também é uma distribuição unimodal, e o grupo modal é o ( 50 ; 60 ], igual ao grupo modal da prova, como um todo, e acima do grupo modal para a Formação Geral.
Na Tabela 6.6, é apresentada uma comparação dos resultados em relação à Categoria Administrativa, à Organização Acadêmica e à Modalidade de Ensino, agora levando -se em conta o desempenho de alunos do Componente de Conhecimento Específico da prova.
No que se refere à Organização Acadêmica, a maior Média foi a dos Centros Universitários ( 59,0 ), vindo a seguir a das Universidades ( 56,9 ), a dos CEFET/IFET ( 50,0 ) e, por fim, a das Faculdades ( 49,3 ).
O maior Desvio padrão, e acima do valor para o Brasil, como um todo, foi o das Faculdades ( 15,5 ).
As Universidades obtiveram a maior nota Máxima ( 90,6 ).
Os Centros Universitários obtiveram nota Máxima 86,4 ; os CEFET/IFET, 81,7 ; e as Faculdades obtiveram nota Máxima 79,7.
A maior Mediana foi encontrada nos Centros Universitários ( 59,0 ) e a menor, nos CEFET/IFET ( 49,3 ), sendo as demais 58,3, nas Universidades, 51,3, nas Faculdades.
A menor nota Mínima foi zero, encontrada em Universidades, e a maior nota Mínima foi 17,8, nos Centros Universitários.
Nas Faculdades e nos CEFET/IFET, as notas mínimas foram, respectivamente, 7,4 e 11,1.
Observa -se que existe diferença estatisticamente significativa ao nível de 95 \% no Componente de Conhecimento Específico entre as notas das Universidades e as notas das Faculdades e as dos CEFET/IFET, e entre as notas dos Centros Universitários e as das Faculdades e as dos CEFET/IFET Quanto à Categoria Administrativa, observa -se um comportamento semelhante àquele da parte de Formação Geral e à prova, como um todo, ou seja, não existe diferença estatisticamente significativa entre as médias das IES Públicas ( 55,9 ) e as das IES Privadas ( 55,1 ).
Neste caso, também, a maior Média foi obtida por alunos de IES Públicas de ensino.
Quanto ao comportamento das médias, no que toca à Modalidade de Ensino, não há diferença estatisticamente significativa entre a Média da Educação Presencial ( 55,5 ) e a da Educação a Distância ( 55,9 ). 6.2 ANÁLISE DAS QUESTÕES OBJETIVAS Esta seção apresenta estatísticas selecionadas e histogramas das Questões Objetivas de Formação Geral ( 6.2.1 ) e de Conhecimento Específico ( 6.2.2 ).
São também apresentadas e comparadas as médias das subpopulações caracterizadas por Grande Região. 6.2.1 Componente de Formação Geral Na Tabela 6.7, são apresentadas as Estatísticas Básicas relativas às oito questões objetivas do componente da prova que abrange a Formação Geral dos estudantes.
A Média do Brasil foi 46,1.
A menor Média foi encontrada na região Centro-Oeste ( 39,4 ), e a maior, nas regiões Nordeste e Sudeste ( 48,3 ).
As demais médias foram : 44,1, na região Norte e 47,4, na região Sul.
O Desvio padrão do Brasil foi 21,9, sendo o maior Desvio padrão encontrado na região Sul ( 23,3 ), e o menor, na região Norte ( 18,0 ).
Os demais desvios foram : 21,1, na região Nordeste, 21,7, na região Sudeste e 20,9, na região Centro-Oeste.
As medianas do Brasil, como um todo, e de três das cinco regiões foram 50,0 – Nordeste, Sudeste e Sul.
A região Norte apresentou mediana igual a 43,8 e a região Centro-Oeste, a menor mediana, igual a 37,5.
A nota Máxima 100,0 foi alcançada em três das Grandes Regiões.
As exceções foram as regiões Norte e Centro-Oeste, com nota Máxima 87,5.
As notas mínimas ( 0,0 ) foram iguais para quatro das cinco regiões, exceto a região Norte, com nota mínima igual a 12,5.
Na Tabela 6.8, são apresentados o Índice de Facilidade e o Índice de Discriminação ( Ponto-Bisserial ) para cada uma das questões objetivas do Componente de Formação Geral.
Quanto ao Índice de Facilidade, foram usadas as seguintes cores para diferenciar o nível de dificuldade da questão :  Azul para as questões classificadas com índice Muito fácil ( \&gt ; =0,86 ), verde para as questões classificadas com índice Fácil ( 0,61 a 0,85 ), amarelo para as questões classificadas com Médio ( 0,41 a 0,60 ), vermelho para as questões classificadas com Difícil ( 0,16 a 0,40 ) e roxo para as questões classificadas com Muito difícil ( \&lt ; =0,15 ).
Já quanto ao Índice de Discriminação, foram usadas as seguintes cores para qualificar a questão :  As questões classificadas com índice Fraco receberam a cor vermelha ( \&lt ; =0,19 ), as classificadas com Médio receberam a cor amarela ( 0,20 a 0,29 ), as classificadas com Bom receberam a cor verde ( 0,30 a 0,39 ) e as classificadas com Muito bom ( \&gt ; =0,40 ) receberam a cor azul.
As questões objetivas do Componente de Formação Geral, segundo o Índice de Facilidade, foram assim avaliadas : das oito questões, nenhuma teve o Índice de Facilidade classificado como Muito fácil, e uma questão foi tida como Fácil, por ter índice de acertos 0,61.
Cinco questões foram consideradas com índice de dificuldade Médio, situando -se no intervalo entre 0,41 e 0,60 do Índice de Facilidade, ou seja, houve entre 41,0 \% e 60,0 \% de acertos, enquanto duas questões foram classificadas na categoria Difícil, situando -se no intervalo entre 0,16 e 0,40.
Nenhuma questão apresentou menos de 15 \% de acertos, razão pela qual não houve questão classificada como Muito difícil.
O Índice de Facilidade variou de 0,27 a 0,61, e o de Discriminação, de 0,34 a 0,51.
Na Tabela 6.9, é apresentada a distribuição das questões, levando -se em conta, simultaneamente, a classificação dos dois índices.
As sete questões com Índice de Discriminação Muito bom figuraram entre três níveis de dificuldade Fácil, Médio e Difícil : uma classificada na categoria Fácil ( a questão 8 ) do Índice de Facilidade, quatro, na categoria Médio ( questões 2, 3, 5 e 7 ), e duas, na categoria Difícil ( questões 1 e 4 ).
Em particular, a questão 5 foi a que apresentou o maior poder discriminatório, com índice 0,51, e foi considerada Médio em termos de facilidade, com uma proporção de 0,41 acertos.
O máximo classificadas com Difícil ( 0,16 a 0,40 ) e roxo para as questões classificadas com Muito difícil ( \&lt ; =0,15 ).
Já quanto ao Índice de Discriminação, foram usadas as seguintes cores para qualificar a questão : • As questões classificadas com índice Fraco receberam a cor vermelha ( \&lt ; =0,19 ), as classificadas com Médio receberam a cor amarela ( 0,20 a 0,29 ), as classificadas com Bom receberam a cor verde ( 0,30 a 0,39 ) e as classificadas com Muito bom ( \&gt ; =0,40 ) receberam a cor azul.
As questões objetivas do Componente de Formação Geral, segundo o Índice de Facilidade, foram assim avaliadas : das oito questões, nenhuma teve o Índice de Facilidade classificado como Muito fácil, e uma questão foi tida como Fácil, por ter índice de acertos 0,61.
Cinco questões foram consideradas com índice de dificuldade Médio, situando -se no intervalo entre 0,41 e 0,60 do Índice de Facilidade, ou seja, houve entre 41,0 \% e 60,0 \% de acertos, enquanto duas questões foram classificadas na categoria Difícil, situando -se no intervalo entre 0,16 e 0,40.
Nenhuma questão apresentou menos de 15 \% de acertos, razão pela qual não houve questão classificada como Muito difícil.
O Índice de Facilidade variou de 0,27 a 0,61, e o de Discriminação, de 0,34 a 0,51.
Na Tabela 6.9, é apresentada a distribuição das questões, levando -se em conta, simultaneamente, a classificação dos dois índices.
As sete questões com Índice de Discriminação Muito bom figuraram entre três níveis de dificuldade Fácil, Médio e Difícil : uma classificada na categoria Fácil ( a questão 8 ) do Índice de Facilidade, quatro, na categoria Médio ( questões 2, 3, 5 e 7 ), e duas, na categoria Difícil ( questões 1 e 4 ).
Em particular, a questão 5 foi a que apresentou o maior poder discriminatório, com índice 0,51, e foi considerada Médio em termos de facilidade, com uma proporção de 0,41 acertos.
O máximo de acertos foi alcançado pela questão 8 com um Índice de Facilidade de 0,61.
As questões 1 e 4 foram consideradas Difícil, com índices de Facilidade 0,40 e 0,27, respectivamente.
A questão de número 6 apresentou um Índice de Facilidade de 0,45 ( Médio ), ou seja, menos da metade dos estudantes conseguiu resolvê -la, dentro do universo de participantes.
Seu Índice de Discriminação foi Bom ( 0,34 ).
No Gráfico 6.4, para exemplificar, analisa -se o comportamento da questão de número 4 de Formação Geral.
Trata -se de uma questão considerada Difícil, em relação à facilidade e a que obteve o segundo maior Índice de Discriminação dessa parte da prova.
Neste gráfico, cada uma das cinco curvas representa o percentual de respostas em determinada alternativa da questão, em função do número de acertos dos estudantes nessa parte da prova ( Formação Geral/Múltipla Escolha ), antes de possíveis eliminações pelo critério Ponto-Bisserial.
A curva em verde corresponde à alternativa B, a alternativa correta para essa questão.
Observa -se que entre os estudantes com menor número de acertos, nessa parte do exame, a situação mais frequente foi a escolha da alternativa D ( em laranja ), incorreta.
Por exemplo, dentre os estudantes que acertaram três questões, 31,4 \% escolheram a alternativa C ( em preto ), 23,3 \% escolheram a alternativa E ( em vermelho ), 18,4 \% escolheram a alternativa D ( em laranja ), 15,7 \% escolheram a alternativa A ( azul ) e 11,2 \% escolheram a alternativa B ( em verde ).
Dentre os que acertaram três respostas entre as questões de múltipla escolha de Formação Geral, nenhum deixou a questão em branco, e nenhum marcou mais de uma alternativa, invalidando a questão.
À medida que o número de acertos aumenta, indicando desempenho melhor nessa parte da prova, aumenta, concomitantemente, a proporção de estudantes que selecionaram a alternativa correta B, atingindo 100 \% para os estudantes com oito acertos.
Essa análise permite verificar como a questão discriminou os grupos de desempenho, justificando -se o alto índice obtido na questão.
Cumpre notar que não é possível inferir deste gráfico nem o índice de facilidade ( que seria uma média da proporção ponderada pela quantidade de alunos com cada uma das notas ), nem o índice de discriminação Ponto-Bisserial, por razão equivalente.
No caso extremo, no qual a grande concentração dos acertos dos alunos fosse abaixo de quatro, o índice de facilidade seria obrigatoriamente abaixo de 25 \% ( neste exemplo ).
Caso a concentração fosse em seis acertos ou mais, o índice seria obrigatoriamente acima de 60 \%.
Os gráficos relativos às demais questões de Formação Geral constam do Anexo I. 6.2.2 Componente de Conhecimento Específico Na Tabela 6.10, são apresentadas as Estatísticas Básicas em relação às questões objetivas do Componente de Conhecimento Específico da prova de Secretariado Executivo por Grande Região.
Nessa parte da prova, duas questões foram anuladas pela Comissão Assessora da Área, as questões 15 e 31.
Além disso, como será discutido ainda nesta seção, duas questões objetivas do Componente de Conhecimentos Específicos não foram usadas no cômputo das notas, por terem sido descartadas pelo critério Ponto-Bisserial.
Assim, as notas foram calculadas com base em 23 das 27 questões objetivas de Conhecimentos Específicos.
A Média do Brasil desse componente foi 55,6.
A menor Média foi observada na região Centro-Oeste ( 50,2 ), e a maior, na região Sudeste ( 58,7 ).
O Desvio padrão de todo o Brasil foi 15,9, sendo o menor Desvio padrão encontrado na região Sudeste ( 14,4 ), e o maior, na região Centro-Oeste ( 17,2 ).
A Mediana de todo o Brasil foi 56,5, o mesmo valor da Mediana encontrada nas regiões Nordeste e Sul.
Nas regiões Norte e Centro-Oeste, a Mediana foi menor : 52,2, em cada uma dessas regiões, e, na região Sudeste, a Mediana foi 60,9.
A nota Máxima da prova foi 95,7, obtida nas questões objetivas do Componente de Conhecimento Específico, por, pelo menos, um aluno da região Nordeste.
A nota Mínima foi zero, obtida na região Centro-Oeste.
Nas demais regiões, as notas mínimas foram : 21,7, na Norte, 8,7, no Nordeste e no Sul, e 4,3, no Sudeste.
Na Tabela 6.11, são apresentados os Índices de Facilidade e Discriminação ( Ponto-Bisserial ) das questões objetivas do Componente de Conhecimento Específico, para os estudantes de Secretariado Executivo.
Para facilitar a diferenciação das questões, serão usadas as mesmas cores da Tabela 6.8 para as diferentes classificações dos Índices de Facilidade e de Discriminação.
Observa -se que duas questões foram anuladas pela CAA e que as análises dizem respeito às 25 questões válidas.
A partir do Índice de Facilidade obtido, pode -se concluir que mais de dois quintos das questões objetivas da prova foram considerados, pelo menos, Médio ( 11 questões ).
Das 25 questões válidas, cinco foram classificadas como Difícil ou como Muito difícil.
Nove questões foram classificadas como Fácil, e nenhuma como Muito Fácil.
Já quanto aos índices de discriminação das questões objetivas do Componente de Conhecimento Específico da prova, obtém -se como resultado a seguinte classificação : 15 das 25 questões foram consideradas boas, enquanto três delas tiveram Índice de Discriminação notas foram calculadas com base em 23 das 27 questões objetivas de Conhecimentos Específicos.
A Média do Brasil desse componente foi 55,6.
A menor Média foi observada na região Centro-Oeste ( 50,2 ), e a maior, na região Sudeste ( 58,7 ).
O Desvio padrão de todo o Brasil foi 15,9, sendo o menor Desvio padrão encontrado na região Sudeste ( 14,4 ), e o maior, na região Centro-Oeste ( 17,2 ).
A Mediana de todo o Brasil foi 56,5, o mesmo valor da Mediana encontrada nas regiões Nordeste e Sul.
Nas regiões Norte e Centro-Oeste, a Mediana foi menor : 52,2, em cada uma dessas regiões, e, na região Sudeste, a Mediana foi 60,9.
A nota Máxima da prova foi 95,7, obtida nas questões objetivas do Componente de Conhecimento Específico, por, pelo menos, um aluno da região Nordeste.
A nota Mínima foi zero, obtida na região Centro-Oeste.
Nas demais regiões, as notas mínimas foram : 21,7, na Norte, 8,7, no Nordeste e no Sul, e 4,3, no Sudeste.
Na Tabela 6.11, são apresentados os Índices de Facilidade e Discriminação ( Ponto-Bisserial ) das questões objetivas do Componente de Conhecimento Específico, para os estudantes de Secretariado Executivo.
Para facilitar a diferenciação das questões, serão usadas as mesmas cores da Tabela 6.8 para as diferentes classificações dos Índices de Facilidade e de Discriminação.
Observa -se que duas questões foram anuladas pela CAA e que as análises dizem respeito às 25 questões válidas.
A partir do Índice de Facilidade obtido, pode -se concluir que mais de dois quintos das questões objetivas da prova foram considerados, pelo menos, Médio ( 11 questões ).
Das 25 questões válidas, cinco foram classificadas como Difícil ou como Muito difícil.
Nove questões foram classificadas como Fácil, e nenhuma como Muito Fácil.
Já quanto aos índices de discriminação das questões objetivas do Componente de Conhecimento Específico da prova, obtém -se como resultado a seguinte classificação : 15 das 25 questões foram consideradas boas, enquanto três delas tiveram Índice de Discriminação Muito bom.
Assim, para quase três quartos das questões – 18 em 26 – os Índices de Discriminação foram Bom ou Muito bom.
Dentre as demais, cinco delas foram classificadas como Médio, e outras duas, como Fraco, sendo sete, por conseguinte, a quantidade de questões nos dois patamares mais baixos de discriminação.
Constata -se, assim, que a prova – no que se refere ao Componente de Conhecimento Específico – possuía boa capacidade de discriminar dentre aqueles que dominam ou não o conteúdo.
Na Tabela 6.12, é apresentada a distribuição das questões levando -se em conta, simultaneamente, a classificação dos dois índices.
Dentre as questões que alcançaram os maiores índices de discriminação, uma delas foi classificada com Índice de Discriminação Muito bom, a de número 21.
Seu índice foi 0,44, e, quanto ao índice de Facilidade, foi classificada como Fácil, com 78 \% dos estudantes marcando a opção correta.
A questão de número 22 foi a mais difícil dentre as 25 questões específicas válidas, com baixo Índice de Muito bom.
Assim, para quase três quartos das questões – 18 em 26 – os Índices de Discriminação foram Bom ou Muito bom.
Dentre as demais, cinco delas foram classificadas como Médio, e outras duas, como Fraco, sendo sete, por conseguinte, a quantidade de questões nos dois patamares mais baixos de discriminação.
Constata -se, assim, que a prova – no que se refere ao Componente de Conhecimento Específico – possuía boa capacidade de discriminar dentre aqueles que dominam ou não o conteúdo.
Na Tabela 6.12, é apresentada a distribuição das questões levando -se em conta, simultaneamente, a classificação dos dois índices.
Dentre as questões que alcançaram os maiores índices de discriminação, uma delas foi classificada com Índice de Discriminação Muito bom, a de número 21.
Seu índice foi 0,44, e, quanto ao índice de Facilidade, foi classificada como Fácil, com 78 \% dos estudantes marcando a opção correta.
A questão de número 22 foi a mais difícil dentre as 25 questões específicas válidas, com baixo Índice de Facilidade, apenas 11 \% de acertos.
Essa questão apresentou poder discriminatório muito baixo, 0,14.
Destaca -se, também, a questão 26, com Índice de Facilidade 0,23, o que, em termos percentuais, corresponde a 23 \% de estudantes que responderam acertadamente e cujo Índice de Discriminação foi 0,03.
Tais questões foram, portanto, pelo critério Ponto-Bisserial, consideradas inadequadas.
Por isso, as questões 22 e 26 foram eliminadas do cômputo da nota final.
A título de exemplo das análises do comportamento das questões objetivas, no Gráfico 6.5, analisa -se a questão 21 do Componente de Conhecimento Específico.
Essa questão foi considerada, pelas respostas dos estudantes avaliados na prova, como uma questão com nível de facilidade Fácil, 0,78, ou seja, 78 \% dos estudantes assinalaram, acertadamente, a opção C, correspondente ao gabarito.
Como já comentado, seu Índice de Discriminação foi igual a 0,44, classificado como Muito bom, o maior valor de discriminação.
Neste gráfico, cada uma das cinco curvas representa o percentual de respostas em determinada alternativa da questão 21, em função do número de acertos dos estudantes nessa parte da prova, antes de possíveis eliminações de questões pelo critério Ponto-Bisserial.
A alternativa correta C, representada no gráfico pela curva em preto, foi escolhida em maiores proporções pelos alunos com desempenho melhor nessa parte da prova.
Já as alternativas incorretas, também denominadas distratores, foram selecionadas, principalmente, por aqueles com notas mais baixas.
Observa -se que a soma não é 100 \%, por causa das questões não respondidas ou com mais de uma opção marcada.
Aqueles com nota zero, na sua quase totalidade, deixaram essa questão em branco ou marcaram mais de uma alternativa, comportamento considerado inválido.
A proporção de alunos que selecionaram a resposta correta C aumenta gradativamente, chegando a atingir 100 \%, a partir de 20 acertos ou mais, enquanto a proporção dos que escolheram alternativas incorretas decai, a partir de cinco ou seis acertos, em função do número de acertos nessa parte da prova.
Facilidade, apenas 11 \% de acertos.
Essa questão apresentou poder discriminatório muito baixo, 0,14.
Destaca -se, também, a questão 26, com Índice de Facilidade 0,23, o que, em termos percentuais, corresponde a 23 \% de estudantes que responderam acertadamente e cujo Índice de Discriminação foi 0,03.
Tais questões foram, portanto, pelo critério Ponto-Bisserial, consideradas inadequadas.
Por isso, as questões 22 e 26 foram eliminadas do cômputo da nota final.
A título de exemplo das análises do comportamento das questões objetivas, no Gráfico 6.5, analisa -se a questão 21 do Componente de Conhecimento Específico.
Essa questão foi considerada, pelas respostas dos estudantes avaliados na prova, como uma questão com nível de facilidade Fácil, 0,78, ou seja, 78 \% dos estudantes assinalaram, acertadamente, a opção C, correspondente ao gabarito.
Como já comentado, seu Índice de Discriminação foi igual a 0,44, classificado como Muito bom, o maior valor de discriminação.
Neste gráfico, cada uma das cinco curvas representa o percentual de respostas em determinada alternativa da questão 21, em função do número de acertos dos estudantes nessa parte da prova, antes de possíveis eliminações de questões pelo critério Ponto-Bisserial.
A alternativa correta C, representada no gráfico pela curva em preto, foi escolhida em maiores proporções pelos alunos com desempenho melhor nessa parte da prova.
Já as alternativas incorretas, também denominadas distratores, foram selecionadas, principalmente, por aqueles com notas mais baixas.
Observa -se que a soma não é 100 \%, por causa das questões não respondidas ou com mais de uma opção marcada.
Aqueles com nota zero, na sua quase totalidade, deixaram essa questão em branco ou marcaram mais de uma alternativa, comportamento considerado inválido.
A proporção de alunos que selecionaram a resposta correta C aumenta gradativamente, chegando a atingir 100 \%, a partir de 20 acertos ou mais, enquanto a proporção dos que escolheram alternativas incorretas decai, a partir de cinco ou seis acertos, em função do número de acertos nessa parte da prova.
Os gráficos relativos às demais questões do Conhecimento Específico constam do Anexo I. 6.3 ANÁLISE DAS QUESTÕES DISCURSIVAS Esta seção apresenta estatísticas selecionadas e histogramas das Questões Discursivas de Formação Geral ( 6.3.1 ) e Conhecimento Específico ( 6.3.2 ).
São também apresentadas e comparadas as médias de subpopulações, caracterizadas por Grande Região. 6.3.1 Componente de Formação Geral As análises dos resultados de desempenho dos estudantes de Secretariado Executivo nas duas questões discursivas relativas à Formação Geral encontram -se na Tabela 6.13 e no Gráfico 6.6.
Na Tabela 6.13, observa -se que a nota Média nesse conjunto de questões foi abaixo da obtida nas objetivas.
Os estudantes de todo Brasil obtiveram, em Formação Geral, Média 46,1, nas questões objetivas e 36,9, nas questões discursivas.
Pode -se notar, também, que o Desvio padrão nesse conjunto de questões foi menor do que o obtido nas objetivas : 21,8, nas questões objetivas e 19,7, nas questões discursivas.
A maior Média foi obtida na região Sul ( 40,2 ), e a menor, na região Centro-Oeste ( 29,1 ).
A Mediana de todo o Brasil, neste componente, foi 39,5.
Nas regiões Nordeste, Sudeste e Sul foram maiores, respectivamente, 39,8, 43,5 e 43,0.
Já nas regiões Norte e Centro-Oeste, as medianas foram menores, respectivamente, 32,0 e 30,0.
A nota Máxima ( 89,0 ) foi obtida na região Nordeste, sendo 70,5 a Máxima na região Norte ; 80,5, na Sudeste ; 81,5, na Sul ; e 75,5, na região Centro-Oeste.
A nota Mínima ( 0,0 ) foi a mesma em todas as regiões do Brasil.
No Gráfico 6.6, está representada a distribuição das notas nas questões discursivas do Componente de Formação Geral.
A moda dessa distribuição ocorre no intervalo ( 40 ; 50 ], com frequência de 22,6 \%, seguida pelo intervalo ( 30 ; 40 ], com frequência de 17,8 \%.
Destaca -se também, como um máximo local, o intervalo [ 0 ; 10 ], com distribuição em torno de 13,4 \% do total de notas, sendo que, no intervalo [ 0 ; 10 ], inclui -se, além da nota zero, a frequência de alunos que deixaram esse tipo de questão em branco.
Na sequência, os resultados verificados para cada uma das questões discursivas de Formação Geral serão apresentados, estabelecendo -se relações com os conteúdos abordados em cada uma delas.
Os comentários da Banca de docentes corretores, a respeito do observado na correção das respostas dos estudantes, suas impressões e conclusões serão apresentados junto à análise de cada questão.
Cumpre esclarecer que, tendo -se em vista que as questões discursivas de Formação Geral são padronizadas, ou seja, constam de todas as provas, os comentários da Banca são os mesmos para todas as carreiras acadêmicas, sendo direcionados a todos os estudantes que participaram do Enade/2018.
A seguir, serão analisados os desempenhos da Área de Secretariado Executivo nas duas questões discursivas de Formação Geral do Enade/2018, comparando -se os resultados obtidos com comentários para cada questão. 6.3.1.1 Análise de Conteúdo da Questão Discursiva 1 do Componente de Formação Geral Os dados de Secretariado Executivo, obtidos a partir das respostas à questão 1, encontram -se na Tabela 6.14 e no Gráfico 6.7.
Nessa questão – a de melhor desempenho entre as duas de Formação Geral – os alunos de todo o Brasil tiveram Média 36,8.
A maior Média para a questão 1 foi obtida na região Sudeste ( 39,7 ), e a menor, na região Centro-Oeste ( 29,5 ).
Quanto à variabilidade das notas, o Desvio padrão de todo o Brasil foi 24,5.
O menor Desvio padrão foi obtido na região Norte ( 21,8 ), e o maior Desvio padrão foi obtido na região Nordeste ( 25,2 ).
Apenas a região Centro-Oeste apresentou mediana 30,0.
Nas outras quatro regiões e no Brasil como um todo, a mediana foi 40,0.
A nota máxima da questão discursiva 1 foi 100,0, obtida na região Nordeste.
Nas demais, as máximas foram : 75,0, na região Norte, 90,0, nas regiões Sudeste e Sul, e 80,0, na região Centro-Oeste.
Já a nota mínima para o Brasil foi zero, nota mínima obtida em todas as regiões.
No Gráfico 6.7, mostra -se a distribuição das notas na questão discursiva 1 do Componente de Formação Geral.
Observa -se que a maior frequência corresponde ao intervalo ( 50 ; 60 ], seguida pelo intervalo ( 40 ; 50 ].
Já o intervalo dos alunos que deixaram a questão em branco corresponde a 14,9 \% dos participantes. 6.3.1.2 Comentários sobre a correção de Conteúdo das respostas à Questão Discursiva 1 A Questão 1 exigia do aluno um posicionamento sobre conteúdo humanista e crítico.
O comando da questão solicitava a redação de um texto que formulasse e articulasse argumentos consistentes envolvendo temas previstos na Portaria nº 444, de 30 de maio de 2018, que dispõe sobre o componente de Formação geral do Enade 2018, tais como : solidariedade/violência, tolerância/intolerância, inclusão/exclusão, sexualidade, relações de gênero e relações étnico-raciais ( art. 7º, item VIII ), bem como ética, democracia e cidadania ( art. 7º, item I ).
O enunciado da questão buscava o estabelecimento da relação de três dimensões – os Direitos Humanos, a Democracia e as Fake News.
Para isso, recorreu a quatro suportes : a reprodução de uma foto e três fragmentos adaptados de textos que foram disponibilizados em páginas da Internet ( da Anistia Internacional, do Jornal El País-Brasil e do Labic-Laboratório de estudos sobre Imagem e Cibercultura ).
Os fragmentos traziam dados sobre o aumento do número de assassinatos de grupos minoritários, as origens e o trabalho da vereadora Marielle Franco, sua luta pelos Direitos Humanos, seu assassinato, e a produção de fake news.
O comando solicitava um texto abordando dois tópicos : o tensionamento entre a defesa dos Direitos Humanos e a produção de notícias falsas, tendo como base a vereadora Marielle – a defesa dos Direitos Humanos realizada por ela e as notícias falsas divulgadas após a sua morte.
As fake news também estavam presentes no segundo tópico de análise : os prejuízos que provocam para a sociedade democrática.
Por um lado, os textos selecionados para constituir o enunciado da questão continham apoio suficiente para formulação e articulação de argumentos para discorrer sobre o tensionamento entre a atuação da Marielle em defesa dos Direitos Humanos e fake news.
Por outro, havia ênfase nos aspectos biográficos e na trajetória da vereadora Marielle, o que levou a respostas pouco focadas no comando, que discutiam a vida e a morte da vereadora, sua importância como um ícone dos Direitos Humanos, opiniões em torno do seu assassinato, até mesmo sobre a responsabilidade pelo crime, muitos textos com a presença/defesa de fake news ou produção de novas.
O foco sobre Marielle também ensejou respostas basicamente centradas na questão da violência.
Foram frequentes as interpretações com viés político-partidário, não previstas no padrão de resposta.
O item do comando que solicitava a discussão sobre democracia e fake news encontrou menos apoio nos textos.
Apesar de se esperar que graduandos tivessem maturidade, inclusive política, para demonstrar alguma compreensão dos aspectos principais que envolvem o convívio democrático, esse aspecto foi pouco discutido e, quando presente, quase sempre foi apenas mencionado, numa quase reprodução do comando.
A falta de entendimento do que é o conceito de Direitos Humanos também foi percebida nas respostas, seja para recomendá -lo como importante para a sociedade, seja para criticá -lo como algo que defende bandidos.
Foi marcante a ideia que vincula Direitos Humanos a direitos apenas das minorias.
Não foi incomum encontrar respostas que personificavam os Direitos Humanos, tratando -os como sujeito de alguma ação que deveria ser implementada na sociedade em geral ou em casos específicos.
As fake news foram citadas muitas vezes, mas pouco explicadas, conforme previsto no padrão de resposta.
O argumento utilizado geralmente vinculava as consequências de fake news para a imagem da Marielle, ou seja, centrava no fato em si, já indicado no enunciado, sem qualquer análise ou aprofundamento próprio.
É importante salientar que muitas respostas, provavelmente de grupos vinculados a cursos específicos, continham argumentações recorrentes : os Direitos Humanos como algo vinculado à Constituição e a tratados dos quais o Brasil é signatário ; as fake news como produto do mundo digital e globalizado, relacionado à era da pós-verdade ; a discussão sobre o papel das mídias sociais e da imprensa. produto do mundo digital e globalizado, relacionado à era da pós-verdade ; a discussão sobre o papel das mídias sociais e da imprensa.
Como já apontado, houve baixa quantidade de respostas contemplando o segundo tópico do comando – prejuízos da produção de notícias falsas para a sociedade democrática ; fato motivado, possivelmente, pela falta de um elemento que ressaltasse esse ponto nos textos motivadores.
De modo geral, os estudantes não conseguiram recorrer a conhecimentos tácitos sobre democracia para construir argumentos próprios.
As respostas que abordavam o tema, geralmente estavam vinculadas à ideia do direito à verdade, do direito à integridade da pessoa ou dos direitos gerais garantidos pela Constituição.
A relação da questão com o momento político brasileiro foi indicada em um número expressivo de respostas.
Por vezes, apenas questionando a temática da questão – caso Marielle – em uma prova para avaliar o ensino superior no Brasil ou a importância dos Direitos Humanos na atual conjuntura política do país, levando à anulação da resposta.
Houve também interpretações das consequências do uso das fake news no processo eleitoral de 2018, influenciando a escolha consciente e livre das pessoas.
As respostas demonstraram atenção político social frente a esses temas, porém com pouco domínio conceitual do campo de conhecimento sobre Direitos Humanos, fake news e democracia.
Em termos da linguagem adotada pelos alunos, há dois grandes grupos de respostas.
No primeiro grupo, os que responderam de forma clara, com autonomia de pensamento e explicitação de justificativas para os aspectos solicitados no enunciado, demonstrando conhecimento sobre o que significam Direitos Humanos, fake news e democracia, além da articulação entre essas questões.
Por outro lado, um segundo grupo dos que não conseguiram ir além do que foi citado no enunciado, apenas repetindo -o, muitas vezes desviando o foco da questão para temas correlatos, como a vida de Marielle ou a violência.
Nesse segundo grupo, mais facilmente foram encontradas respostas desconexas e mal articuladas, com textos confusos que explicitam o não conhecimento sobre o tema e a falta de adequação à linguagem, necessária a um concluinte do ensino superior.
As abordagens dos Direitos Humanos mais recorrentes nas respostas corretas foram : a referência da Marielle como representante/ícone dos DH ; e os DH como direitos previstos na Constituição e em tratados internacionais dos quais o Brasil é signatário.
No entanto, também foram encontradas críticas aos Direitos Humanos, tais como : DH vistos como “ direitos de bandidos ” ; a defesa pouco embasada de frases de efeito como “ direitos humanos para humanos direitos ”.
Também chamou atenção o tratamento dos DH como um personagem, ou um órgão público, evidenciando desconhecimento sobre o tema.
Em relação às fake news, dois argumentos previstos pelo padrão de respostas foram os mais presentes : a noção de que o mundo globalizado e a facilidade de reprodução de informações pelas mídias atuais levam à criação de mentiras que são compartilhadas e que geram consequências ; a Marielle como alvo de fake news que tentaram desconstruir sua luta e seu legado.
Também foi frequente a ideia de que a produção de informações falsas é realizada por uma imprensa tendenciosa com intenção de manipular a informação.
No item relacionado a fake news e democracia, as respostas corretas afirmavam que as fake news inviabilizam o acesso à verdade e à garantia da integridade de um Estado Democrático de Direito.
Também foi abordada a influência das fake news nos resultados das eleições dos últimos anos, nos Estados Unidos e no Brasil, gerando consequências para a vida das pessoas.
Por fim, um outro aspecto abordado nas respostas a este tópico foi a polarização da sociedade em termos políticos, com a afirmação de que as fake news dificultam o debate, estimulam posições extremas e rivalidades prejudiciais à convivência em uma democracia plena.
A notas mais fracas ( de zero a 35 ) foram cerca de 20 \% das respostas corrigidas – as que não estavam em branco nem foram desconsideradas ou anuladas.
Foram respostas que se restringiam a cópia ou paráfrases do enunciado ; as que tinham foco em desdobramentos políticos e policiais do caso Marielle ; as que tinham argumentação confusa e/ou incompleta.
A grande maioria das notas, em torno de 65 \%, foi mediana, entre 40 a 70.
Nesse grupo ficaram as respostas que continham alguma abordagem adequada sobre Direitos Humanos, sobre as atividades da vereadora Marielle e sobre as fake news divulgadas após sua morte.
As notas mais altas deste grupo foram para respostas que incluíram a questão do tensionamento, quase sempre, em um texto coerente e consistente.
Muitas vezes, a nota mediana foi consequência do baixo índice de respostas à segunda parte do comando, que solicitava a discussão sobre os prejuízos das fake news para a democracia.
As melhores notas só foram obtidas por aqueles que incluíram a discussão sobre os prejuízos das fake news para a democracia.
Apenas 15 \% das respostas corrigidas receberam avaliação igual ou superior a 75.
A deficiência principal detectada pelos corretores foi a dificuldade de grande parte dos alunos para interpretar o enunciado da questão e para se expressar em textos consistentes.
De modo geral, os temas foram tratados de forma fragmentada e pouco analítica.
Além disso, chamou à atenção a predominância do lugar comum na construção argumentativa das respostas, a superficialidade do debate em torno da democracia e o desconhecimento de conceitos como Direitos Humanos, esses confundidos com instituições e até pessoas. 6.3.1.3 Análise de Conteúdo da Questão Discursiva 2 do Componente de Formação Geral A Tabela 6.15 mostra que o desempenho médio dos estudantes na questão discursiva 2 ( média 29,4 ) foi inferior ao obtido na questão discursiva 1 ( média 36,8 ).
A região Sul foi aquela cuja Média foi maior ( 33,8 ), enquanto a de menor Média, assim como na questão discursiva 1, foi a região Centro-Oeste ( 20,5 ).
Quanto à variabilidade das notas, o Desvio padrão de todo o Brasil foi 22,0, inferior ao obtido na questão discursiva 1 ( 24,5 ).
O maior desvio nessa questão foi obtido na região Norte ( 23,4 ), enquanto o menor foi obtido na região Centro-Oeste ( 19,7 ).
A Mediana de todo o Brasil foi 25,0, a mesma da região Nordeste.
Nas regiões Norte e Centro-Oeste, as medianas foram 20,0, e nas regiões Sudeste e Sul, 35,0.
A nota máxima no Brasil foi 100,0, a mesma das regiões Nordeste e Sul, sendo 85,0, nas demais regiões.
Já a nota mínima foi zero, mesma de todas as regiões.
No Gráfico 6.8, mostra -se a distribuição das notas na questão discursiva 2 do Componente de Formação Geral.
Observa -se que a maior frequência corresponde ao intervalo ( 30 ; 40 ], com cerca de um quinto dos participantes, seguido dos intervalos ( 40 ; 50 ], com frequência de 15,4 \%, e dos alunos que deixaram a questão em branco, 13,9 \%.
Nota -se, ainda, que as notas ficaram menos dispersas em comparação às da questão discursiva de número 1, o que pode ser constatado, também, pela comparação do desvio padrão das notas da questão discursiva 2 ( 22,0 ) e o das notas da questão discursiva 1 ( 24,5 ). 6.3.1.4 Comentários sobre a correção de Conteúdo das respostas à Questão Discursiva 2 Os textos e o comando da Questão 2 de Formação Geral são perfeitamente adequados ao que se espera avaliar em Formação Geral, ao tratar das funções sociais dos museus, contextualizando a temática sob o ponto de vista da memória e da perspectiva de futuro.
Na questão, estavam presentes, principalmente, o conteúdo humanista e crítico, a interpretação e a produção de textos sobre temas como cultura e arte, educação, ciência e democracia, previstos na Portaria que dispõe sobre o componente de Formação Geral do Enade 2018.
Os conhecimentos exigidos para responder ao comando, explicitados no padrão de resposta, vão além do senso comum, por se exigir a compreensão de conceitos como “ função social ” ( de instituições sociais ), “ memória ” ( no sentido de memória social ), “ perspectiva de futuro ” ( no contexto sócio-histórico-cultural ) e “ sociedade contemporânea ” ( que deveria independer da interpretação de cada grupo social a respeito da sociedade na qual vive ).
No entanto, apesar de envolver tais conceitos, espera -se que alunos concluintes de cursos superiores de graduação sejam capazes de discutir funções de instituições sociais, especialmente porque o enunciado trazia fato recente e de grande repercussão como apoio.
O padrão de resposta elencou as funções sociais que poderiam ser abordadas nas respostas dos estudantes, como também argumentos que poderiam ser usados ao longo das mesmas.
A tarefa dos estudantes era a construção de um texto expositivo que articulasse a importância de museus, suas funções sociais e as perspectivas de memória e de futuro, o que se mostrou de dificuldade de média a difícil.
Apesar de o enunciado ser claro, muitos estudantes, provavelmente influenciados pelos textos, limitaram -se a discutir aspectos relacionados ao Museu Nacional e, em particular, ao incêndio que o destruiu.
Foram poucos os que interpretaram corretamente o escopo do comando, que solicitava a discussão das funções sociais de museus de uma forma geral, no sentido de argumentar a importância deste tipo de instituição.
As quatro funções sociais elencadas no padrão de resposta foram encontradas na quase totalidade das redações dos estudantes de forma articulada com a importância dos museus.
Entretanto, quase nunca os textos abordavam mais do que uma função social e vinculavam, de forma clara, a importância desse tipo de instituição para a sociedade contemporânea.
Além disso, a solicitação de contextualização sob o foco museu/memória/perspectiva de futuro parece não ter sido bem compreendida e, em consequência, esteve pouco presente nas respostas.
Também se observou confusão entre os conceitos de história e memória e uma visão estreita do significado de sociedade contemporânea, às vezes, restrita à cidade do respondente, outras, ao Rio de Janeiro, ou ainda a um tipo de sociedade idealizada.
A relevância que os textos de apoio deram ao caso particular do incêndio do Museu Nacional do Rio de Janeiro e aos prejuízos causados ao seu acervo, induziu a maioria dos estudantes a elaborar respostas tratando das funções a partir da perda e, em particular, à perda do tipo de acervo daquele museu e do tipo de pesquisa que lá se desenvolve, divergindo do comando da questão, que solicitava texto abordando funções sociais de museus em geral.
Também foram encontradas muitas respostas que se limitavam a discutir o que o estudante considerava como possíveis causas do incêndio, a falta de manutenção, questões políticas e de gestão, sem responder ao que foi solicitado.
Foi também recorrente a transcrição de frases dos textos do enunciado nas respostas, diretamente ou por paráfrases, revelando certa dificuldade de criação textual.
A análise das respostas dos estudantes, além de revelar dificuldade de abordar a temática esperada, revelou outros dois aspectos predominantes.
Um deles, como já comentado, foi a ênfase no incêndio do Museu Nacional, que orientou grande parte dos textos dos estudantes.
Outro, também evidenciando desvio do esperado, foi o foco na perda patrimonial.
Nesse grupo, houve respostas que até apresentaram uma função social, mas olhada pelo ponto de vista da perda.
Nessa linha, muitas respostas enalteciam a importância do Museu Nacional, sua relevância no cenário nacional, o impacto da perda de seu acervo, isso complementado por frágeis reflexões sobre as funções sociais.
Em consequência, a perspectiva de desdobramentos mais ricos, que tratassem de aspectos sociais relacionados a diferentes tipos de museu foi prejudicada.
Muitas respostas trataram da preservação da memória e/ou guarda do acervo ; identidade, tanto individual quanto coletiva ; conhecimento sobre antepassados ; evolução da sociedade e dos seus costumes ; valorização e disseminação do conhecimento por meio da liberação do acesso aos acervos ; bem como atividades pedagógicas e exposições de arte.
O aspecto produção de pesquisa e conhecimento, embora frequente, foi mencionado de forma muito simples, sem qualquer aprofundamento, para além do que já era mencionado no enunciado.
Frequentemente apresentava -se o museu como produtor de pesquisas para o campo da saúde, o que evidencia desconhecimento das reais possibilidades de pesquisa, mesmo no Museu Nacional.
Quanto a perspectiva de futuro, quando presente no texto, era mencionada de forma breve, com afirmações genéricas sobre a possibilidade de melhoria do futuro pelo simples conhecimento do passado, numa percepção de que o passado determina o futuro, de forma linear.
Foram raras as respostas com alguma argumentação que associasse pesquisa, memória, identidade e reflexão social como base para a construção de alternativas de futuro.
A relação entre turismo e museu apareceu considerando a visita a museus como forma de lazer e, algumas vezes, como geradora de recursos.
O desempenho pode ser considerado fraco.
Muitas respostas deixavam claro que o estudante nunca tinha visitado um museu e, de uma forma geral, observou -se desconhecimento sobre suas funções sociais.
Não se esperava que história, memória social e museologia fossem assuntos dominados pelos estudantes das diversas áreas que participaram do Enade 2018, mas esperava -se que essa temática pudesse ser abordada do ponto de vista das experiências pessoais ou de conhecimento social.
Como uma vivência própria com os museus também se revelou frágil ou inexistente, o conhecimento sobre o assunto ficou bem restrito aos textos motivadores e ao senso comum, construído, basicamente, pelo noticiário em torno do incêndio.
O universo dos argumentos dos estudantes se mostrou muito limitado.
Cada dimensão foi praticamente representada por uma ideia simplista.
De forma recorrente : a preservação da memória para construção da identidade de um povo ; a compreensão do passado por meio das pesquisas como oportunidade para melhorar o futuro ; o museu como um espaço em que há um aprendizado estimulado pelo lazer.
Houve, também, uma grande concentração de respostas associando o museu a um local de guarda de acervo e de história.
As respostas à Questão 2 evidenciaram dificuldades conceituais.
Além dos conceitos de história, memória e museu se misturarem, parecendo não se distinguirem, ao mencionarem a pesquisa, também se observou falta de distinção entre pesquisas colegiais, com objetivos de aprendizado, e as realizadas por pesquisadores, com objetivos de avançar cientificamente.
A falta de compreensão do que se pedia no comando da questão e a dificuldade de expressão linguística foram dois aspectos marcantes na avaliação dos corretores.
Foi frequente os textos apresentarem ideias soltas ou de forma confusa, sem constituir unidade.
O vocabulário usado se revelou bem restrito, e fortemente apoiado no coloquial, na linguagem oral.
Predominou o senso comum em lugar de uma reflexão mais formal envolvendo os conceitos envolvidos. 6.3.1.5 Análise de Língua Portuguesa das Questões Discursivas do Componente de Formação Geral Os dados de Secretariado Executivo, obtidos a partir das respostas às questões discursivas do Componente de Formação Geral, no que tange à Língua Portuguesa, encontram -se na Tabela 6.16 e no Gráfico 6.9.
Nesse aspecto, os alunos de todo o Brasil obtiveram Média 52,4.
A maior Média em relação à Língua Portuguesa foi obtida na região Sudeste ( 56,1 ), e a menor, na região Centro-Oeste ( 45,4 ).
Quanto à variabilidade das notas, o Desvio padrão de todo o Brasil foi 22,7.
O menor Desvio padrão foi obtido na região Sudeste ( 20,5 ) e o maior Desvio padrão foi obtido na região Centro-Oeste ( 24,8 ).
A Mediana das notas de Língua Portuguesa foi 60,0 para o Brasil, como um todo, a mesma obtida na região Nordeste.
A Mediana da região Sul foi 65,0, a maior dentre as obtidas pelas Grandes Regiões, enquanto a da Centro-Oeste foi 55,0, a menor.
Na região Norte, a Mediana foi 56,3, e na região Sudeste, 62,5.
A nota Máxima para todo o Brasil foi 85,0, com, pelo menos, um aluno tirando essa nota nas regiões Nordeste, Sudeste e Centro-Oeste.
Na região Norte, a nota Máxima foi 75,0, e na região Sul, 82,5.
Já a nota Mínima foi zero em todas as regiões do país.
No Gráfico 6.9, mostra -se a distribuição das notas de Língua Portuguesa do Componente de Formação Geral.
Observa -se que a maior frequência ( 37,4 \% ) corresponde à dos alunos que obtiveram nota no intervalo ( 60 ; 70 ].
Destacam -se, também, os alunos que obtiveram nota no intervalo ( 50 ; 60 ], correspondendo a 23,0 \% dos participantes.
Já os alunos que deixaram ambas as questões em branco representam quase 10 \% do total. 6.3.1.6 Comentários sobre a correção das respostas de Formação Geral com respeito à Língua Portuguesa As duas questões discursivas de Formação Geral ensejaram duas configurações textuais distintas, tendo em vista os seus encaminhamentos : a Questão 1 teve um caráter mais opinativo e a Questão 2, mais expositivo, conforme se explicita a seguir.
Ambas as questões permitiram a observação do desempenho linguístico dos participantes em situação formal, competência exigida no processo de formação universitária.
As habilidades de produção textual envolvidas têm, igualmente, como condição básica, o domínio do padrão formal culto da Língua Portuguesa, no que diz respeito aos aspectos relativos às questões de caráter ortográfico, textual, morfossintático e vocabular.
O encaminhamento das questões evidencia a expectativa de que o participante utilizasse seus conhecimentos sobre os dois temas e estruturasse seus textos de acordo com as características do registro formal adequado à situação comunicativa.
Tal configuração determina exigências quanto : à adequação da seleção vocabular, ao desenvolvimento do conteúdo, à estruturação sintática dos períodos, à organização lógica das ideias, à utilização de procedimentos de encadeamento textual e de referenciação, à obediência às exigências morfossintáticas próprias da modalidade escrita da norma-padrão, ao respeito às regras ortográficas e às regras de acentuação gráfica.
Esse encaminhamento exigiu do graduando o domínio do modo de organização textual de base dissertativo-argumentativa, essencial ao processo de formação universitária.
As habilidades de produção textual envolvidas têm como condição básica o domínio do padrão formal culto da Língua Portuguesa, no que diz respeito aos aspectos relativos às questões de caráter ortográfico, textual, morfossintático e vocabular.
O enunciado da Questão 1 tomava como base três fragmentos de textos jornalísticos publicados em sites da internet, que abordavam o aumento do assassinato de minorias e de defensores de direitos humanos no Brasil, entre eles a vereadora Marielle Franco.
O comando solicitava que o participante desenvolvesse um texto dissertativo e argumentativo com as temáticas Fake News, Democracia e Direitos Humanos.
Embora não tenha havido referência explícita à tipologia textual esperada para o desenvolvimento da resposta, essa proposta encaminhou o participante para a elaboração de um texto opinativo, com o objetivo de discutir o problema e manifestar um posicionamento crítico.
O enunciado da Questão 2 abordava o incêndio do Museu Nacional do Rio de Janeiro, ocorrido em 2018, e seus impactos na ciência brasileira e internacional, tomando como base dois fragmentos de textos jornalísticos publicados em sites da internet.
O comando solicitava que o participante desenvolvesse um texto abordando três aspectos da função social dos museus na sociedade contemporânea, encaminhando -o para a elaboração de um texto expositivo que confirmasse a posição defendida nos fragmentos motivadores a respeito da importância dessas instituições.
Para atender a essas exigências, o padrão de resposta utilizado na avaliação das questões discursivas de Formação Geral considerou aspectos relevantes ao bom desempenho linguístico como competências distintas, de modo a permitir um mapeamento detalhado do domínio dos recursos disponíveis na Língua Portuguesa para a comunicação escrita formal : ortográficos, textuais, morfossintáticos e vocabulares ( ver detalhes no Anexo VIII – Padrão de Respostas ).
Aspectos Ortográficos Para avaliar a competência relativa ao domínio das convenções ortográficas da norma-padrão da Língua Portuguesa, observou -se o respeito às regras de acentuação gráfica e da grafia padrão das palavras ( com ausência de abreviaturas próprias da linguagem da internet ), de acordo com as convenções estabelecidas pela legislação em vigor e consubstanciadas no Vocabulário Ortográfico da Língua Portuguesa, editado pela Academia Brasileira de Letras.
Aspectos Textuais Para avaliar a competência relativa ao domínio dos procedimentos de estruturação textual, procurou -se observar a utilização de mecanismos linguístico-discursivos que estabeleçam o encadeamento lógico entre as partes do texto, de forma a garantir a progressão e a coerência textuais.
Esta competência destaca -se por sua relevância, já que são avaliadas a clareza e a unidade textual.
Aspectos Morfossintáticos Para avaliar o domínio dos diferentes aspectos morfossintáticos próprios da modalidade escrita formal da norma-padrão da Língua Portuguesa, observou -se se o texto produzido atendeu às seguintes exigências : a concordância nominal, a concordância verbal, a regência nominal, a regência verbal, a flexão nominal, a flexão verbal, a correlação entre os tempos verbais, a colocação pronominal e a utilização de sinais de pontuação que contribuam para a organização lógica da frase e do texto.
Aspectos Vocabulares Para avaliar a adequação da seleção vocabular à modalidade escrita formal da Língua Portuguesa, exigida pela situação comunicativa, observou -se o atendimento às seguintes exigências : precisão na seleção/utilização do vocabulário relacionado à temática solicitada pela questão ; ausência de marcas da oralidade, como termos de sentido muito genérico ( “ coisa ”, “ negócio ”, “ você ” ) e termos de registros mais informais ( como gírias, jargões, frases feitas, ditados populares, termos regionais ).
Assim, espera -se que o participante respeite a adequação vocabular não usando gírias ou expressões coloquiais, evite repetição desnecessária de palavras e utilize um vocabulário mais formal, como solicitado por um texto dissertativo.
Para efeito de pontuação, a avaliação do desempenho linguístico considerou três grandes grupos de competências, agrupando aspectos explicitados, com os seguintes pesos relativos : aspectos ortográficos ( 20 \% ) ; aspectos textuais ( 40 \% ), aspectos morfossintáticos e vocabulares ( 40 \% ).
Os critérios de pontuação utilizados no padrão de resposta procuraram atender à diversidade de desempenho dos graduandos.
Nos textos analisados, observaram -se, assim, vários níveis de atendimento às exigências do padrão formal próprio da modalidade escrita da Língua Portuguesa, configurando um"continuum".
Por um lado, textos que traduzem com exatidão a expectativa de uso formal da língua, em todos os níveis analisados, marcados por clareza e unidade ; e, no outro extremo, textos quase incompreensíveis devido ao grande número de desvios de todos os tipos, principalmente os de caráter textual.
As duas questões discursivas de Formação Geral analisadas revelaram desempenhos distintos dos participantes.
A Questão 1 propiciou a oportunidade de maior desenvolvimento do tema solicitado, gerando consequências na estruturação textual, já que os textos foram mais longos, mais elaborados e fluentes.
Em virtude da polêmica sobre a existência das fake news e sobre a visão pública da personagem da vereadora Marielle Franco, observaram -se respostas que polemizaram o tema, favorecendo a construção argumentativa do texto.
A Questão 2 teve um comportamento distinto, devido ao caráter mais expositivo do enunciado e ao reduzido conhecimento dos participantes sobre o tema.
Assim, os textos foram mais curtos, frequentemente com repetição de ideias dos textos motivadores.
Quanto aos aspectos linguísticos analisados durante esta avaliação, observaram -se os seguintes resultados : Aspectos ortográficos : O desempenho dos participantes revelou uma diferença muito grande nos dois aspectos analisados nesta competência : baixo índice de desvios da grafia padrão e grande índice de desvios de acentuação.
Em vários casos, ocorre ausência completa de acentuação gráfica em todas as palavras do texto.
Os resultados revelaram que a tendência dominante entre os universitários brasileiros é a eliminação da acentuação gráfica, talvez motivada pelos hábitos relacionados às redes sociais e pela ausência de esclarecimento dos meios de comunicação, das autoridades e das escolas sobre as decisões do Acordo Ortográfico da Língua Portuguesa de 1990.
Os casos mais sistemáticos de eliminação do acento indicador da sílaba tônica foram : - palavras proparoxítonas ( por exemplo “ generos ”, “ arbitro ”, “ politico ”, “ publicas ”, “ numero ” ) ; - palavras paroxítonas terminadas em ditongo crescente ( por exemplo “ area ”, “ noticiarios ”, “ individuo ”, “ historia ”, “ varios ”, “ noticia ”, “ negligencia ”, “ ocorrencia ” ) ; - palavras oxítonas ( por exemplo “ ninguem ”, “ esta ”, “ ate ”, “ tambem ” ).
Por outro lado, destaca -se o uso indevido do acento gráfico em palavras como “ melâncolia ” e “ intervênção ”, por exemplo.
Quanto ao domínio das convenções relativas à grafia das palavras, observaram -se desvios esporádicos como : hipercorreção pela escolha de “ e ” no lugar de “ i ”, por influência de hábitos da oralidade ( “ descriminação ” no lugar de “ discriminação ”, “ entervenção ” no lugar de “ intervenção ” ) ; desvios de grafia relacionados à variação diastrática, como por exemplo “ perca ”, “ subjulgadas ”, “ precoseito ”, “ soubre ”, “ vecendo ”, “ indesencia ”, “ apolojia ”, “ fulga ”, “ dereitos ” ;...
Observaram -se, também, muitos casos de inadequação no uso da maiúscula : uso indevido para destacar determinadas palavras-chave do texto, como “ Brasileiros ”, “ Homens ” ; ausência de maiúscula para grafar nomes próprios, como"marielle".
Expressões como"Direitos Humanos"foram grafadas diversificadamente, com maiúsculas e com minúsculas.
Destaque -se, também, a maioria dos participantes que grafam os textos inteiramente em caixa alta, o que impede a identificação de desvios relacionados ao uso de maiúsculas.
Vale observar, também, que, ao contrário do que se esperava, não apareceram abreviaturas próprias do “ internetês ”, relacionadas ao uso de redes sociais e e-mails, como por exemplo"vc".
Aspectos textuais : Para alguns dos graduandos essa competência se revelou como a mais problemática, tendo em vista os inúmeros problemas observados, desvios acumulados durante toda a formação escolar.
São eles : sequência justaposta de ideias sem encaixamentos sintáticos, provocando uma redução de estruturas subordinadas, ao lado do aumento na frequência de estruturas coordenadas e absolutas ; frases fragmentadas que comprometem a estrutura lógico-gramatical ; frases formadas apenas por oração subordinada, sem oração principal ; redução no uso de conectores para expressar relações lógicas essenciais à construção do texto como consequência da mudança de estruturação frasal ; emprego inadequado do pronome relativo ( com omissão da preposição ou a utilização do pronome relativo inadequado, como “ onde ” ), refletindo hábitos da oralidade ; ausência de recursos de referenciação, como substituição de termos por sinônimos, hiperônimos, nominalizações, expressões metafóricas. lógico-gramatical ; frases fragmentadas que comprometem a estrutura lógico-gramatical ; frases formadas apenas por oração subordinada, sem oração principal ; redução no uso de conectores para expressar relações lógicas essenciais à construção do texto como consequência da mudança de estruturação frasal ; emprego inadequado do pronome relativo ( com omissão da preposição ou a utilização do pronome relativo inadequado, como “ onde ” ), refletindo hábitos da oralidade ; ausência de recursos de referenciação, como substituição de termos por sinônimos, hiperônimos, nominalizações, expressões metafóricas.
Esses problemas revelam dificuldades em relação à estrutura formal do texto produzido, o que é preocupante ao se levar em conta que são graduandos em fase final de formação.
Em alguns textos, falta um mínimo de textualidade e de domínio do registro padrão da língua.
Na verdade, observaram -se relações linguísticas quase agramaticais, como as estabelecidas pela sequência de gerúndios sem o apoio de um ponto de partida para a organização das informações gramaticais e semânticas.
Quanto à utilização dos sinais de pontuação, observou -se uma ausência dos mesmos na maioria dos textos analisados.
Foi muito frequente a ocorrência de parágrafos sem marca interna de pontuação para separar os períodos.
São os seguintes os tipos de problemas mais encontrados :  vírgula : utilização de vírgula para separar o sujeito e o predicado ; uso de vírgula no lugar do ponto para separar ideias que constituem períodos distintos ; ausência de vírgula para separar oração adjetiva explicativa ;  ponto e vírgula : utilização do ponto e vírgula no lugar de vírgula ;  ponto final : ausência de ponto final para separar períodos.
Aspectos morfossintáticos e vocabulares : Os resultados são transparentes em relação aos aspectos mais problemáticos no desempenho dos participantes nestas duas competências.
Em relação aos aspectos morfossintáticos, seguem algumas observações sobre os desvios mais frequentes.
Quanto à regência, o desvio mais frequente foi a falta do sinal indicativo da crase – isso revela que o usuário não tem consciência de que, sob a forma do termo “ a ”, existe a presença de uma contração entre a preposição “ a ” ( exigida pela regência do termo anterior ) e o artigo definido “ a ”.
Outro problema relacionado à regência verbal e à nominal, encontrado frequentemente nas respostas, foi a ausência de preposição antes de pronome relativo, processo generalizado na modalidade oral da língua, em situações de registro informal.
Apesar da possibilidade de que essa alteração de regência se generalize no padrão escrito da Língua Portuguesa, como já está ocorrendo até em textos jornalísticos, o não emprego da preposição foi considerado como desvio neste processo de avaliação.
Outro desvio de regência significativo foi a utilização inadequada de uma preposição ou sua ausência após o verbo ou o nome ( substantivo ou adjetivo ).
A concordância verbal e a concordância nominal apresentaram alguns desvios muito frequentes.
Quanto à concordância de número, observou -se, como apontam as pesquisas nessa área, ausência de marca ( com sujeito anteposto ou posposto ) ou uso indevido ( uso inadequado da marca de plural comandado pelo núcleo plural da locução adjetiva, apesar de o substantivo que funciona como núcleo do sintagma nominal estar no singular ).
Uma ocorrência que se destacou foi a ausência de acento circunflexo na forma plural do presente do indicativo do verbo “ ter ”, que foi considerada como um desvio na concordância verbal e não na acentuação gráfica.
Quanto à concordância de gênero, vários casos foram observados, normalmente no âmbito de sintagmas nominais longos, em que o adjetivo está afastado do substantivo.
Quanto à questão da colocação pronominal, foram poucos os desvios observados.
Concluiu -se que, no registro escrito formal, a maioria dos participantes já incorporou regras como a não introdução da frase por um pronome oblíquo e a próclise na presença de um termo atrator.
Várias marcas de oralidade foram identificadas, embora não com alta frequência : uso do pronome relativo “ onde ” como relativo universal, falta de artigo definido antes de substantivo isolado ou antes de termos em enumerações.
Quanto aos aspectos vocabulares, alguns tipos de inadequação foram observados : expressões da oralidade ; seleção vocabular incompatível com o contexto, gerando falta de inteligibilidade ; falta de domínio de vocabulário mais abstrato e de maior complexidade, essencial ao desenvolvimento do texto de base dissertativa.
O principal aspecto observado foi a excessiva repetição de certas palavras, como o termo"pessoa", por exemplo, revelando limitação de repertório vocabular.
Em função do tema solicitado na Questão 1, diferentes grafias da expressão fake news não foram consideradas.
Para sintetizar, os problemas que mais se destacaram nos textos analisados foram : a ) truncamentos sintáticos, caracterizados por estruturas frasais incompreensíveis devido à ruptura da complexidade sintática própria da modalidade escrita ; na modalidade oral da língua, em situações de registro informal.
Apesar da possibilidade de que essa alteração de regência se generalize no padrão escrito da Língua Portuguesa, como já está ocorrendo até em textos jornalísticos, o não emprego da preposição foi considerado como desvio neste processo de avaliação.
Outro desvio de regência significativo foi a utilização inadequada de uma preposição ou sua ausência após o verbo ou o nome ( substantivo ou adjetivo ).
A concordância verbal e a concordância nominal apresentaram alguns desvios muito frequentes.
Quanto à concordância de número, observou -se, como apontam as pesquisas nessa área, ausência de marca ( com sujeito anteposto ou posposto ) ou uso indevido ( uso inadequado da marca de plural comandado pelo núcleo plural da locução adjetiva, apesar de o substantivo que funciona como núcleo do sintagma nominal estar no singular ).
Uma ocorrência que se destacou foi a ausência de acento circunflexo na forma plural do presente do indicativo do verbo “ ter ”, que foi considerada como um desvio na concordância verbal e não na acentuação gráfica.
Quanto à concordância de gênero, vários casos foram observados, normalmente no âmbito de sintagmas nominais longos, em que o adjetivo está afastado do substantivo.
Quanto à questão da colocação pronominal, foram poucos os desvios observados.
Concluiu -se que, no registro escrito formal, a maioria dos participantes já incorporou regras como a não introdução da frase por um pronome oblíquo e a próclise na presença de um termo atrator.
Várias marcas de oralidade foram identificadas, embora não com alta frequência : uso do pronome relativo “ onde ” como relativo universal, falta de artigo definido antes de substantivo isolado ou antes de termos em enumerações.
Quanto aos aspectos vocabulares, alguns tipos de inadequação foram observados : expressões da oralidade ; seleção vocabular incompatível com o contexto, gerando falta de inteligibilidade ; falta de domínio de vocabulário mais abstrato e de maior complexidade, essencial ao desenvolvimento do texto de base dissertativa.
O principal aspecto observado foi a excessiva repetição de certas palavras, como o termo"pessoa", por exemplo, revelando limitação de repertório vocabular.
Em função do tema solicitado na Questão 1, diferentes grafias da expressão fake news não foram consideradas.
Para sintetizar, os problemas que mais se destacaram nos textos analisados foram : a ) truncamentos sintáticos, caracterizados por estruturas frasais incompreensíveis devido à ruptura da complexidade sintática própria da modalidade escrita ; b ) redução de estruturas subordinadas, compensada pelo aumento na frequência de estruturas coordenadas e absolutas, por um lado, ou pela elaboração de estruturas textuais prejudicadas pelo excesso de ideias sem a devida conexão subordinativa ; c ) redução no uso de conectores para expressar relações lógicas essenciais à construção do texto, substituídos pela exigência de inferência por parte do interlocutor para suprir a sua ausência ; d ) redução no uso do subjuntivo, ao lado da ampliação do uso do indicativo, combinado a estruturas frasais coordenadas ou absolutas ; d ) empobrecimento do processo de referenciação, com a repetição de palavras, sem processos mais sofisticados de substituição ; e ) simplificação da marcação da categoria tempo na morfologia verbal ; f ) desrespeito às regras de concordância verbal e nominal próprias do padrão formal ; g ) ausência do acento grave indicador da crase, revelando falta de conhecimento dos critérios de regência verbal e nominal ; h ) falta de domínio de vocabulário mais abstrato e de maior complexidade, essencial ao desenvolvimento do processo dissertativo ; i ) redução do emprego da acentuação gráfica.
Como resultado da correção, merece atenção a constatação de que a modalidade escrita tem apresentado a tendência a uma evidente simplificação, aproximando -se das características da modalidade oral da Língua Portuguesa.
No caso dos textos de base dissertativa ( expositiva ) e de base argumentativa, inscritos em um registro formal, a distância entre as duas modalidades é ainda maior, o que provoca desvios recorrentes em todos os quatro aspectos analisados : ortográficos, textuais, morfossintáticos e vocabulares.
Destaca -se, nessa avaliação, o comprometimento textual, por meio de estruturas fragmentadas e/ou truncadas, rompendo a complexidade sintática esperada no padrão formal. 6.3.2 Componente de Conhecimento Específico Na parte da prova relativa às questões discursivas do Componente de Conhecimento Específico ( Tabela 6.17 ), observa -se que a Média foi mais alta do que para as questões discursivas do Componente de Formação Geral.
Enquanto no Componente de Formação Geral a Média para estudantes de Secretariado Executivo de todo o Brasil foi 36,9, na parte de Conhecimento Específico, a Média foi 54,8.
A maior Média desse componente foi obtida pelos estudantes da região Sul ( 62,5 ), e a menor, pelos da região Norte ( 38,7 ).
Quanto à variabilidade das notas, o Desvio padrão de todo o Brasil foi 24,1.
O maior Desvio padrão foi encontrado na região Centro-Oeste ( 26,0 ), e o menor, na região Norte ( 19,4 ).
A maior nota Máxima, 100,0, foi obtida nas regiões Sudeste e Sul.
Nas demais regiões, as notas Máximas foram : 80,0, na região Norte, 86,7, na região Nordeste, e 91,7, na região Centro-Oeste.
A nota Mínima ( 0,0 ) foi obtida por, pelo menos, um aluno em todas as regiões do Brasil.
A Mediana do Brasil, como um todo, foi 60,0 ou seja, pelo menos metade dos participantes tirou nota igual ou inferior a 60,0 nas questões discursivas do Componente de Conhecimento Específico.
A maior Mediana foi obtida na região Sul ( 66,7 ), e a menor, na região Norte ( 40,9 ).
No Gráfico 6.10, representa -se a distribuição das notas nas questões discursivas no Componente de Conhecimento Específico.
A moda dessa distribuição ocorre no intervalo [ 60 ; 70 ], com quase 21 \% do total de participantes.
A análise de cada uma dessas questões será feita a seguir. 6.3.2.1 Análise da Questão Discursiva 3 do Componente de Conhecimento Específico Na questão 3, cujos resultados aferidos encontram -se descritos na Tabela 6.18, a Média dos estudantes de todo o Brasil foi 56,0, inferior a da questão 4 e superior a da questão 5.
A menor Média nessa questão foi obtida pelos alunos da região Norte ( 39,7 ), enquanto a maior Média foi obtida na região Sul ( 68,9 ).
Quanto à variabilidade das notas, o Desvio padrão de todo o Brasil foi 31,5.
O maior Desvio padrão foi obtido na região Centro-Oeste ( 33,8 ), enquanto o menor foi obtido na região Sudeste ( 27,7 ).
A nota Máxima, 100,0 pontos, foi alcançada por, pelo menos, um aluno de todas as regiões, com exceção da região Norte que teve nota Máxima 85,0.
A Mediana do Brasil, como um todo, foi zero, a mesma para todas as regiões.
No Gráfico 6.11, mostra -se a distribuição das notas na questão discursiva 3, do Componente de Conhecimento Específico da área de Secretariado Executivo.
Na questão 3 10,3 \% dos estudantes deixaram a resposta em branco e 5,7 \% receberam nota zero.
A moda da distribuição está no intervalo ( 80 ; 90 ], com frequência de 20,2 \%, seguido do intervalo ( 60 ; 70 ], com frequência de 19,0 \%. 6.3.2.2 - Comentários sobre as respostas à Questão Discursiva 3 Na questão discursiva 3, abordou -se um conteúdo importante, previsto nas Diretrizes Curriculares Nacionais, sendo escolhida uma temática que faz parte dos tópicos previstos nos PPC - Projetos Pedagógicos dos Cursos, o que é específica para a formação do profissional em Secretariado Executivo, no que tange à Assessoria – Gestão e Execução de Rotinas Secretariais.
A questão foi apresentada de forma clara e objetiva, com a introdução do conceito básico e com o questionamento bem definido.
Uma vez que, além da apresentação de alguns benefícios, foi solicitada uma explicação desses benefícios, e tal ação envolveu uma reflexão que nem sempre se mostrou clara nas respostas encontradas, a questão pode ser considerada de grau médio de complexidade.
A utilização de técnicas secretariais com tecnologias inovadoras necessária para a atribuição de eficiência e eficácia nos processos, é uma das competências previstas na formação do profissional de Secretaria Executivo.
Além dessa, deve -se mencionar a capacidade de ter um raciocínio lógico, crítico e analítico, que opere com valores e que estabeleça relações formais e causais entre fenômenos e situações organizacionais.
O aluno comprometido com a construção e com a articulação dos conhecimentos científicos, acadêmicos, tecnológicos e estratégicos específicos da sua área, com visão generalista e intersetorial da organização, é capaz de responder adequadamente essa a questão.
A maioria dos estudantes compreendeu a questão, sendo que a variação das notas ocorreu em função : a ) do número de exemplos citados, uma vez que a questão exigia cinco benefícios que a tecnologia oferece, se aplicada à execução das técnicas secretariais com arquivos, para torná -la mais eficiente.
Dessa forma, o não atendimento a esse número de benefícios já interferia na nota aferida.
Especial cuidado foi necessário, pois, em alguns momentos, o estudante mencionou o mesmo benefício, com outras palavras, o que levou à não pontuação de um deles ; b ) da necessidade de explicar os benefícios, e não, de apenas citá -los.
A ausência da explicação impactou na nota final.
A resposta observada com maior frequência foi aquela em que se apresentavam e se explicavam, pelo menos, três ( dos cinco ) benefícios do uso da tecnologia quando aplicada à execução das técnicas secretariais com arquivos, para torná -la mais eficiente.
Muitos alunos optaram por apresentar a resposta em tópicos e, como o comando da questão não definia o formato do texto, todos os formatos foram aceitos.
Foi comum apresentarem um texto subjetivo, com vocabulário pessoal, poucas vezes técnico, embora tangenciando a resposta correta. 6.3.2.3 - Análise da Questão Discursiva 4 do Componente de Conhecimento Específico Na Tabela 6.19, constam as informações relativas à questão 4 do conjunto de questões do Componente de Conhecimento Específico, sendo a questão de melhor desempenho dentre as três discursivas de Conhecimentos Específicos.
A Média geral do Brasil foi 62,4, sendo a maior Média registrada na região Sul ( 67,7 ), e a menor, na região Norte ( 50,8 ).
A nota Máxima ( 100,0 ) foi atingida por, pelo menos, um concluinte de todas as cinco Grandes Regiões.
A Mediana em todo o Brasil foi 70,0, o mesmo valor foi obtido nas regiões Nordeste e Sudeste.
Na demais regiões, as medianas foram : 65,0, na Norte, 80,0, na região Sul, e 60,0 na região Centro-Oeste.
Também foi zero a nota Mínima de todas as regiões.
No Gráfico 6.12, está representada a distribuição de notas da questão discursiva 4, do Componente de Conhecimento Específico.
Essa distribuição tem moda no intervalo ( 70 ; 80 ], com 22,0 \% do total de participantes.
A categoria “ em branco ” constitui um máximo local, com 9,4 \% do total de participantes optando por não responder à questão 4. 6.3.2.4 Comentários sobre as respostas à Questão Discursiva 4 Na questão discursiva 4, abordou -se um conteúdo importante previsto nas Diretrizes Curriculares Nacionais, sendo escolhida uma temática que faz parte dos tópicos previstos nos PPC - Projetos Pedagógicos dos Cursos –, e que é específica específico para a formação do profissional em Secretariado Executivo, no que tange ao tópico de Assessoria – Redação Empresarial e Oficial em Língua Portuguesa e Inglesa.
A questão foi apresentada de forma clara e objetiva, com a introdução do conceito básico e com questionamento bem definido.
A questão pode ser considerada de grau baixo de complexidade, uma vez que é uma ação frequente da prática do profissional de Secretariado Executivo.
Para uma resposta, adequada o aluno deveria utilizar recursos de Comunicação e Expressão compatíveis com o exercício profissional, inclusive nos processos de negociação e nas comunicações interpessoais ou intergrupais.
Tal reflexão é competência prevista na formação, quando se prevê que ele seja eficaz na utilização de técnicas secretariais, com renovadas tecnologias, imprimindo segurança, credibilidade e fidelidade no fluxo de informações, atendendo ao perfil de um profissional consciente das implicações e das responsabilidades éticas do seu exercício profissional.
A produção de e-mail deveria ser de domínio do profissional em questão, já que requer conhecimentos básicos, que fazem parte das habilidades necessárias a um bom profissional.
Apesar de a linguagem do enunciado ter sido clara e direta, apresentando todos os elementos necessários para a elaboração do e-mail, muitos estudantes não inseriram todos os elementos solicitados e necessários.
Essa questão foi a que apresentou melhor desempenho dos estudantes quando comparada às demais discursivas.
A utilização do espaço destinado à resposta apresentou variação, pois muitos buscaram criar espaços para apresentar uma similaridade visual ao e-mail.
A maioria dos estudantes obedeceu à delimitação de espaço oferecida na prova e desenvolveram os textos a partir da primeira linha.
Poucos estudantes iniciaram suas respostas a partir da segunda linha ou mais.
O padrão de resposta contemplou todos os elementos necessários à redação do e-mail, e o critério de pontuação foi compatível com a solicitação.
Não é possível apresentar um padrão único que garanta todas as possibilidades de redação, pois para cada item necessário no texto, várias possibilidades puderam ser apresentadas, mas a forma como foi apresentado o padrão de resposta permitiu o entendimento e ofereceu flexibilidade para que o corretor o analisasse dentro dos parâmetros.
Ficou claro o grau de facilidade da questão, uma vez que todos os aspectos exigidos para a composição do e-mail foram apresentados no enunciado.
O único caso não previsto era o relativo à ‘ f ’ — assinatura/identificação do remetente e como este fato não foi abordado, as assinaturas foram aceitas.
A maioria dos estudantes produziu adequadamente o e-mail, ocorrendo variação da nota em função da apresentação de todos os requisitos solicitados no enunciado.
Era necessário que o e-mail contemplasse informações solicitadas no enunciado, isentas de julgamento, uma vez que não cabia ao remetente fazer juízo de valor. 6.3.2.5 Análise da Questão Discursiva 5 do Componente de Conhecimento Específico Na Tabela 6.20, constam as informações relativas à questão 5 do conjunto do Componente de Conhecimento Específico.
O desempenho dos estudantes de todo o Brasil nesta questão foi o mais baixo dentre as três questões discursivas desse componente.
A nota Média dos estudantes de todo o Brasil foi 45,8.
A maior Média foi registrada na região Sul ( 51,1 ), enquanto a menor Média foi registrada na região Norte ( 25,6 ).
Quanto à variabilidade das notas, o Desvio padrão dos alunos do Brasil, como um todo, foi 29,3.
O Gráfico 6.13 apresenta a distribuição das notas da questão discursiva 5 do Componente de Conhecimento Específico.
A moda é o intervalo ( 40 ; 50 ], com 23,2 \% dos participantes.
Destacam -se dois máximos locais – a categoria “ em branco ”, com quase 14 \% dos estudantes tendo deixado a resposta em branco, e o intervalo ( 70 ; 80 ], com as notas de 16,7 \% dos participantes.
Enquanto o maior desvio foi encontrado na região Sul ( 30,8 ), o menor foi encontrado na região Norte ( 22,5 ).
A Mediana para o Brasil e em todas as regiões foi 50,0, a mesma das regiões Nordeste, Sudeste e Sul.
Nas regiões Norte e Centro-Oeste, as medianas foram, respectivamente, 30,0 e 40,0.
Para o conjunto de alunos de Secretariado Executivo do Brasil, a nota Máxima foi 100,0, obtida em todas as regiões, exceto na Norte, cuja nota Máxima foi 75,0.
A nota Mínima do Brasil foi zero, o mesmo valor encontrado em todas as regiões.
O Gráfico 6.13 apresenta a distribuição das notas da questão discursiva 5 do Componente de Conhecimento Específico.
A moda é o intervalo ( 40 ; 50 ], com 23,2 \% dos participantes.
Destacam -se dois máximos locais – a categoria “ em branco ”, com quase 14 \% dos estudantes tendo deixado a resposta em branco, e o intervalo ( 70 ; 80 ], com as notas de 16,7 \% dos participantes. 6.3.2.6 Comentários sobre as respostas à Questão Discursiva 5 A questão discursiva 5 abordou um conteúdo importante previsto nas Diretrizes Curriculares Nacionais, escolhendo uma temática que faz parte dos tópicos previstos nos PPC - projetos pedagógicos dos cursos, abordando tema específico para a formação do profissional em Secretariado Executivo, no que tange à Secretariado Executivo estratégica, Relações de poder, cultura e clima organizacional.
A questão foi apresentada de forma clara e objetiva, sem dificuldade de interpretação, mas com muitos elementos a serem discutidos e organizados numa reflexão de, no máximo, 15 linhas.
A questão apresentou uma introdução clara sobre o tema a ser abordado, seguida de um questionamento bem definido, sem possiblidades de outras interpretações, mas com complexidade temática.
Por esse motivo, pode ser considerada de grau difícil de complexidade, pois envolvia uma temática repleta de elementos, propondo a construção de um texto, com abordagem de dois temas correlatos, mas diferentes, da área da gestão da organização, não envolvendo, necessariamente, uma prática frequente no papel do profissional de secretariado executivo, mas necessária quando o mesmo atua em diferentes níveis de gestão da organização.
Tais conceitos precisam estar claramente entendidos para que sejam explanados de forma clara, com a inclusão dos aspectos solicitados, envolvendo, assim, uma reflexão por parte do estudante a qual nem sempre estava presente nas respostas.
Tal reflexão envolveu mais de uma competência prevista na formação do profissional como raciocínio lógico, crítico e analítico, operando com valores e estabelecendo relações formais e causais entre fenômenos e situações organizacionais.
Envolveu, também, habilidade de lidar com modelos inovadores de gestão, domínio dos recursos de expressão e de comunicação compatíveis com o exercício profissional, inclusive nos processos de negociação e nas comunicações interpessoais ou intergrupais, além de receptividade e liderança para o trabalho em equipe, na busca da sinergia.
O aluno, uma vez que está compromissado com a postura crítica e reflexiva articulada com discrição e sensibilidade para as questões transversais do contexto em que atua, tem subsídios para responder à questão.
O padrão de resposta leva em consideração o alinhamento de diversos pontos, como clima organizacional ; relações de poder associados a objetivos e metas organizacionais ; comunicação ; cultura organizacional ; clima organizacional ; e lideranças.
Foi encontrado nas respostas um espectro amplo de abordagens, desde a discussão conceitual dos elementos apresentados na pergunta a uma reflexão adequada dentro dos parâmetros nela propostos.
Foi observado que a maioria dos estudantes, talvez pelo espaço limitado destinado à resposta e tempo reduzido para a elaboração da mesma, em face da complexidade que a questão exigia, optou por abordar, apenas, um dos seguintes aspectos : a ) Relações de Poder e clima organizacional ; b ) liderança e trabalho em equipe.
Os conhecimentos necessários à solução da questão pressupunham o domínio de conceitos quanto à gestão, relações interpessoais e cultura organizacional, assim como a capacidade de articular as temáticas fazendo correlações conforme solicitado no enunciado.
O enfoque escolhido com maior frequência pelos estudantes foi a discussão da liderança e trabalho em equipe, e poucas vezes foi abordado de forma ampla, como solicitado no enunciado da questão.
Poucas respostas apresentaram a correlação poder e clima organizacional, evidenciando o motivo de poucos estudantes apresentarem de forma correta o entendimento do que é o clima organizacional e, consequentemente, a dificuldade em estabelecer sua relação com o poder.
A linguagem utilizada nem sempre foi técnica, o que sugere uma falta de domínio da linguagem específica da área, embora tangenciando a resposta correta.
O espaço destinado à resposta foi insuficiente devido à complexidade da questão.
Para atender à resposta em sua amplitude, o estudante necessitava de um grande poder de síntese, o que em muitos momentos limitou uma articulação adequada do proposto. 6.3.3 Considerações Finais Os alunos que apresentaram boas respostas, em sua maioria, articularam as ideias, apresentaram bons argumentos, souberam conceituar e exemplificar quando necessário.
No entanto, alguns apenas parafrasearam o enunciado, não contemplando a expectativa.

% matched lemmas: 
\end{textsample}
